% !TEX program = xelatex
% ============================================================
%  《电网 AI 数据智能体：原理、方法与实践》—— 绪论
%  编译方式：XeLaTeX 编译一次即可（ctex 自动匹配中文字体）
%  依赖：ctex（TeX Live / MiKTeX 自带），无需外部图片文件
% ============================================================
\documentclass[UTF8,a4paper,12pt]{ctexart}

\usepackage[a4paper,left=2.6cm,right=2.6cm,top=2.8cm,bottom=2.6cm]{geometry}
\usepackage{xcolor}
\usepackage{tikz}
\usetikzlibrary{arrows.meta,positioning,shapes.geometric,fit,backgrounds,calc}
\usepackage{enumitem}
\usepackage{caption}
\usepackage[hidelinks]{hyperref}

% ---------- 配色（深蓝主色 + 红色强调 + 青色/金色辅色）----------
\definecolor{primary}{HTML}{1F4E79}
\definecolor{accent}{HTML}{C00000}
\definecolor{teal}{HTML}{0F8B8D}
\definecolor{gold}{HTML}{B8860B}
\definecolor{lightbg}{HTML}{F2F5F9}

% ---------- 标题样式：彩色分级 ----------
\renewcommand{\thesection}{0.\arabic{section}}
\ctexset{
  section/format  = \Large\bfseries\color{primary},
  section/aftername = \quad,
  section/beforeskip = 1.2em plus .2ex,
  section/afterskip  = 0.7em plus .2ex,
  subsection/format = \large\bfseries\color{teal},
}

% ---------- 图注与编号：图 0-1、图 0-2 …… ----------
\renewcommand{\thefigure}{0-\arabic{figure}}
\captionsetup{font=small,labelsep=quad}

% ---------- 常用 TikZ 样式 ----------
\tikzset{
  bx/.style      = {draw=primary,fill=lightbg,rounded corners=3pt,thick,align=center,inner sep=4pt},
  bxr/.style     = {draw=accent,fill=white,rounded corners=3pt,thick,align=center,inner sep=4pt},
  bxt/.style     = {draw=teal,fill=lightbg,rounded corners=3pt,thick,align=center,inner sep=4pt},
  fl/.style      = {-{Stealth[length=2.4mm]},thick,primary},
}

\begin{document}

% ============================ 标题区 ============================
\begin{center}
  {\Huge\bfseries\color{primary} 绪论}\\[0.35em]
  {\LARGE\color{accent} 从电网数据到数据智能体}\\[0.8em]
  {\small\color{gray} 节号 0.1--0.7 \quad·\quad 配套资源见附录 C、D}
\end{center}

\vspace{0.4em}
\hrule
\vspace{1.0em}

\begin{quote}
\itshape
电网从来不缺数据。真正稀缺的，是把海量数据及时、可靠、可解释地转化为业务判断的能力。
本绪论说明：为什么电网行业需要数据智能体，它与以往的自动化有何不同，本书写给谁、怎么读，
以及全书"六类智能体、一条流水线、多个业务场景"的整体技术路线。
\end{quote}

% ============================ 0.1 ============================
\section{为什么电网行业需要数据智能体}

过去十余年，电网的数字化建设让数据的"存量"发生量级变化。调度侧的 SCADA 与 PMU 以秒级乃至毫秒级
记录潮流与断面，配电侧的 AMI 与用采系统覆盖到台区与用户，运检侧的台账、缺陷与工单不断沉淀，
再叠加 GIS、气象、交易等多源信息，形成了规模庞大、类型复杂的数据资产。然而，"数据多"并不等于
"用得好"。在多数单位，数据增长与业务智能化之间仍横亘着三重矛盾（图~\ref{fig:conflict}）：

\begin{enumerate}[leftmargin=2em,itemsep=2pt]
  \item \textbf{数据体量快速增长，人工分析能力却接近上限。}
        报表口径有限、分析依赖个人经验，难以支撑越来越细的颗粒度与越来越高的时效要求。
  \item \textbf{系统分散、口径不一，业务却需要全局一致的视角。}
        EMS、DMS、营销、GIS 各说各话，同一个"负荷"在不同系统里含义并不相同。
  \item \textbf{专业门槛高，业务对时效的要求却很苛刻。}
        一次异常排查往往要跨越数据、专业算法与业务规则三类知识，非专业人员难以独立完成。
\end{enumerate}

\begin{figure}[htbp]
\centering
\begin{tikzpicture}[font=\small]
\tikzset{
  L/.style={bx,minimum height=10mm,minimum width=4.8cm},
  R/.style={bxr,minimum height=10mm,minimum width=4.8cm}
}
\node[L] (l1) at (0,0)      {数据体量快速增长};
\node[R] (r1) at (6.8,0)    {人工分析能力见顶};
\node at (3.4,0) {\Huge\color{accent}$\neq$};

\node[L] (l2) at (0,-1.7)   {系统分散、口径不一};
\node[R] (r2) at (6.8,-1.7) {业务需要全局一致视角};
\node at (3.4,-1.7) {\Huge\color{accent}$\neq$};

\node[L] (l3) at (0,-3.4)   {专业门槛高};
\node[R] (r3) at (6.8,-3.4) {业务时效要求苛刻};
\node at (3.4,-3.4) {\Huge\color{accent}$\neq$};
\end{tikzpicture}
\caption{电网数据工作的三重矛盾：数据侧在增长，能力侧与业务侧却在原地。}
\label{fig:conflict}
\end{figure}

传统的应对手段各有边界：BI 与报表解决"看得见"，却难以回答"为什么"和"该怎么办"；
规则引擎擅长确定性判断，但规则越写越多、维护越来越难；机器学习模型能做好预测与识别，
却通常只覆盖单点任务，无法串起从取数到报告、再到预警的完整链路。大语言模型与智能体技术的出现，
第一次让"能理解目标、能规划步骤、能调用工具、能自检结果"的\textbf{数据劳动力}成为可能，
这正是本书讨论的起点。

% ============================ 0.2 ============================
\section{从数据工具到数字劳动力}

回顾电网数据工作的手段演进，可以清楚地看到一条"自主性逐级提升"的路径（图~\ref{fig:spectrum}）。
早期是\textbf{工具时代}：Excel、SQL、BI 由人来驱动，工具只负责执行确定动作；随后出现
\textbf{助手（Assistant）}，能被动响应提问，但不主动推进任务；再进一步是\textbf{副驾（Copilot）}，
可以给出建议、与人在同一流程中协作；最终是\textbf{智能体（Agent）}，围绕一个业务目标自主拆解任务、
选择工具、执行并校验结果。

\begin{figure}[htbp]
\centering
\begin{tikzpicture}[font=\small]
\tikzset{bar/.style={bx,minimum height=9mm,align=center}}
\node[bar,minimum width=3.2cm]  (b1) at (0,0)     {助手\\ Assistant\\\footnotesize 被动响应};
\node[bar,minimum width=4.4cm]  (b2) at (5.0,0)   {副驾\\ Copilot\\\footnotesize 协作建议};
\node[bar,minimum width=5.6cm]  (b3) at (10.4,0)  {智能体\\ Agent\\\footnotesize 自主规划与执行};
\draw[fl] (-1.7,-1.15) -- (13.3,-1.15) node[midway,below,black,font=\footnotesize]{自主性逐步增强};
\end{tikzpicture}
\caption{电网数据工作的能力光谱：从被动响应走向面向目标的自主执行。}
\label{fig:spectrum}
\end{figure}

需要强调的是，"数字劳动力"并不等于把已有流程自动化一遍。传统自动化的对象是\textbf{固定流程}，
一旦输入或环境变化就需要人工改脚本；而数据智能体的对象是\textbf{业务目标}，它把"任务怎么拆、
用哪些工具、结果是否可信"作为可推理、可追溯的过程来对待。这一转变，决定了本书讨论的重点
不只是"用哪个模型"，而是"如何把模型、数据、工具与业务规则组织成一条可靠的工作流"。

% ============================ 0.3 ============================
\section{本书的定位}

本书既不是一本 AI 科普读物，也不是一本纯算法专著，而是一本\textbf{面向电网业务、能够照着实践}
的技术读物。它围绕三条线索展开：

\begin{itemize}[leftmargin=2em,itemsep=2pt]
  \item \textbf{原理}——讲清电网数据的特征，以及大模型与智能体的能力边界与物理约束；
  \item \textbf{方法}——给出六类数据智能体的设计方法与工程实现路径；
  \item \textbf{实践}——每章末都落在一个可复现的电网数据任务上，从取数一路做到报告与预警。
\end{itemize}

与市面上偏通用场景的智能体书籍相比，本书的三点不同在于：其一，业务场景全部落在真实的电网
专业语境（调度运行、设备运检、营销客服、新能源与微电网）；其二，始终强调\textbf{可信与可验证}，
把物理约束、确定性计算与人工复核作为强制环节；其三，给出从原型到验收的完整工程路线，
而非停留在演示级样例。

% ============================ 0.4 ============================
\section{读者对象与知识基础}

本书面向两类读者，并尽量降低两端的入门门槛：

\begin{itemize}[leftmargin=2em,itemsep=2pt]
  \item \textbf{电网从业者}——调度、运检、营销、信通与数字化人员，希望借助智能体提高本专业
        的数据工作效率；
  \item \textbf{AI 与软件开发者}——希望把大模型能力真正落到电网业务中的工程师与研究者。
\end{itemize}

在知识基础上，读者只需具备基本的电力常识与基本的编程与数据库直觉即可；对 Transformer、
RAG、工具调用等概念不作前置要求，第 1、2 章会补齐必要的共同知识。

% ============================ 0.5 ============================
\section{本书的三条阅读路径}

不同读者的关注点不同，因此本书提供三条阅读路径（图~\ref{fig:paths}），可各取所需、也可交叉阅读。

\begin{figure}[htbp]
\centering
\begin{tikzpicture}[font=\small]
\tikzset{p/.style={bx,minimum height=12mm,text width=13.4cm,align=left,inner sep=6pt}}
\node[p] (p1) at (0,0)    {\textbf{业务应用路径}：第 1、3 章 $\rightarrow$ 第 5--10 章 $\rightarrow$ 第 13--16 章，先建立概念、再掌握六类能力、最后落到本专业场景。};
\node[p] (p2) at (0,-1.7) {\textbf{技术开发路径}：第 2、4 章 $\rightarrow$ 第 5--12 章 $\rightarrow$ 附录 C、D，先补齐技术底座、再动手实现与编排。};
\node[p] (p3) at (0,-3.4) {\textbf{组织管理路径}：绪论 $\rightarrow$ 第 3 章 $\rightarrow$ 第 17--18 章，关注价值判断、实施路径与风险边界。};
\end{tikzpicture}
\caption{三条阅读路径：业务应用、技术开发与组织管理。}
\label{fig:paths}
\end{figure}

% ============================ 0.6 ============================
\section{本书的整体技术路线}

全书的骨架可以概括为一句话：\textbf{六类智能体、一条流水线、多个业务场景}（图~\ref{fig:route}）。
数据智能体的能力被拆成六类可独立实现、又可串接协作的单元——数据集成、数据治理、数据转换、
洞察、看板与预警；它们在多智能体协作层被编排成一条"采集 $\rightarrow$ 治理 $\rightarrow$ 分析
$\rightarrow$ 报告 $\rightarrow$ 预警"的自主流水线；最后与不同专业的业务模型、知识与流程结合，
形成面向调度、运检、营销与新能源等场景的具体应用。

\begin{figure}[htbp]
\centering
\begin{tikzpicture}[font=\small]
\tikzset{
  ds/.style={bx,minimum height=9mm,align=center,inner sep=4pt},
  ag/.style={bxr,minimum height=11mm,minimum width=2.2cm,align=center,inner sep=3pt},
  pl/.style={bxt,minimum height=10mm,align=center,inner sep=4pt},
  sc/.style={draw=gold,fill=white,rounded corners=3pt,thick,minimum height=9mm,align=center,inner sep=4pt},
  ar/.style={fl}
}
\node[ds,minimum width=14.45cm] (src) at (0,0)
  {数据源：SCADA · PMU · AMI · GIS · 营销 · 气象 · 交易};
\foreach \i/\t in {0/集成,1/治理,2/转换,3/洞察,4/看板,5/预警}{%
  \node[ag] (a\i) at ({(\i-2.5)*2.45},-1.75) {\t};}
\node[pl,minimum width=14.45cm] (pipe) at (0,-3.35)
  {多智能体协作与自主数据流水线（第 11--12 章）};
\node[sc,minimum width=14.45cm] (scene) at (0,-4.95)
  {业务场景：运行调度 · 设备运检 · 营销客服 · 新能源与微电网};
\draw[ar] (0,-0.45) -- (0,-1.18);
\draw[ar] (0,-2.32) -- (0,-2.83);
\draw[ar] (0,-3.87) -- (0,-4.48);
\end{tikzpicture}
\caption{全书技术路线：六类数据智能体、一条自主流水线、多个业务场景。}
\label{fig:route}
\end{figure}

后续各篇即按这一路线展开：第一篇建立共同基础，第二篇界定概念与架构，第三篇逐个剖析六类智能体，
第四篇把它们编排成流水线，第五篇进入专业场景，第六篇讨论工程落地与未来形态。

% ============================ 0.7 ============================
\section{配套实践资源}

本书强调"能照着做"，因此配套资源与正文同等重要，全部收录于附录：

\begin{itemize}[leftmargin=2em,itemsep=2pt]
  \item \textbf{Prompt 模板}（附录 A）——覆盖集成、治理、转换、洞察、看板与预警六类任务；
  \item \textbf{数据模型与字段字典}（附录 B）——统一量测、台区、缺陷、新能源等数据的字段口径；
  \item \textbf{工具与开源框架清单}（附录 C）——选型参考与依赖说明；
  \item \textbf{实践环境与样例数据}（附录 D）——可直接运行的示例程序与完整流水线参考项目；
  \item \textbf{验收与安全检查表}（附录 E）——上线前的自检清单；
  \item \textbf{术语、缩略语与标准索引}（附录 F）——便于跨专业阅读。
\end{itemize}

阅读约定上，本书中的代码以 Python 为主，数据表与流程以图形方式给出；所有结论性分析均要求
给出可核验的计算或证据来源——这既是本书的方法要求，也是电网业务对智能体最基本的要求。

\vspace{1.2em}
\hrule
\vspace{0.6em}
\begin{center}
\small\itshape\color{primary}
下一章，我们从电力系统与电网业务的基础出发，为后续的技术讨论建立共同语言。
\end{center}

\end{document}
